\SourceLeaf{203}\Needspace{10\baselineskip}\begin{center}
[4.]
\end{center}

[Bei \textsc{Lübsen}, \emph{Ausf\"uhrliches Lehrbuch der Elementar-Geometrie}, Hamburg 1851, Seite 11 heisst es:]

\{\textls{Erklärung:} Eine gerade Linie ist diejenige, welche nicht aus ihrer Lage kommt, indem sie sich um ihre beiden festen Endpunkte dreht\HistoricalNote{*)}{\{*)}{So ungef\"ahr h\"orten wir einmal \textsc{Gauss} bei der Erkl\"arung des Fernrohrs und dessen richtigen Gebrauchs den Begriff der geraden Linie festsetzen. Diese Erkl\"arung ist theoretisch fruchtbar, wie die gleich daraus folgenden S\"atze zeigen; ausserdem ist das angegebene Merkmal praktisch wichtig, z.\,B. bei der Justirung eines Fernrohrs, richtigen Bohrung eines Cylinders etc.\}}.\}



\begin{center}
[5.]

\textnormal{THEORIE DES VORTRAGS VON LEHREN, DIE RAUMVERH\"ALTNISSE BETREFFEN.}
\end{center}

1) Alle die Gegenst\"ande, die bei einem Satz relevant sind, \textls{zeichne} man, und bezeichne sie durch Buchstaben oder Ziffern, Punkte durch einen, Linien durch 2, Fl\"achen durch 3, K\"orper allenfalls durch 4, welches jedoch \"ofters auch durch einen geschehen kann.

2) Relativ bewegliche R\"aume, deren Bewegung ohne Stetigkeit bei dem Gesch\"aft in Frage kommt, zeichne man besonders.

3) Die Gr\"ossenrelationen, die relevant sind, bezeichne man auch durch Buchstaben.

4) Man stelle sich das, was man beweisen will, in Form einer Gleichung vor und gehe successive auf deren Gr\"unde zur\"uck.

So wird das ganze Gesch\"aft in seine Elemente gleichsam zerlegt, die man nachher r\"uckw\"arts wieder zusammensetzt. Dabei bemerkt man dann leicht, was beim wirklichen Vortrag zusammengezogen werden kann.

